% Generated draft: requires mathematical and source review.
\documentclass[11pt]{article}
\usepackage[margin=25mm]{geometry}
\usepackage{fontspec}
\setmainfont{Latin Modern Roman}
\usepackage{amsmath,amssymb,mathtools,unicode-math}
\setmathfont{Latin Modern Math}
\usepackage{enumitem,needspace,xurl,hyperref}
\setlength{\parindent}{0pt}
\setlength{\parskip}{5pt}
\setlength{\emergencystretch}{3em}
\newcommand{\N}{\mathbb{N}}
\newcommand{\Z}{\mathbb{Z}}
\newcommand{\Q}{\mathbb{Q}}
\newcommand{\R}{\mathbb{R}}
\newcommand{\C}{\mathbb{C}}
\newcommand{\F}{\mathbb{F}}
\newcommand{\A}{\mathbb{A}}
\newcommand{\PP}{\mathbb{P}}
\newcommand{\E}{\mathbb{E}}
\newcommand{\eps}{\varepsilon}
\newcommand{\dd}{\,\mathrm{d}}
\begin{document}
\section*{UnsolvedMath ID: 30004785}
\textbf{Source problem code:} catalogue entry
\subsection{Definitions and mathematical statement}
Let \(G\) be a split classical group, let \(\widehat G\) denote its complex dual group, and let \(\rho\colon\widehat G\to\operatorname{GL}(V)\) be its standard representation on a complex vector space \(V\).

\textbf{Requested statement.} Prove directly the \((\sigma,\rho)\)-Poisson summation formula for split classical groups \(G\), with \(\rho\) as above.

\textbf{Unresolved formulation.} The supplied record does not define \(\sigma\), specify the base field or the precise family of classical groups, or state the Poisson summation identity. It also does not specify the test-function space, summation domain, transform, normalizations, convergence conditions, or any correction terms required by that identity. Thus the statement above preserves the source's named target, but an explicit identity with complete hypotheses and quantifiers cannot be recovered from the supplied record alone. In particular, no automorphic hypotheses on \(\sigma\) are imposed here by inference. The requirement that the proof be direct is retained, although the record does not define which methods this excludes.

\textbf{Supplied qualification.} The dated curation states that the formula is proved when the relevant global functorial transfer to \(\operatorname{GL}_n\) is known and cuspidal, and that \(\operatorname{GL}_n\) has a direct theorem. It further states that this conditional mechanism does not establish the conjecture for all pairs \((G,\rho)\). These are supplied literature assessments, not independently verified assertions or additional hypotheses completing the missing formulation.
\subsection{Short English statement}
The problem asks for a direct proof of the named \((\sigma,\rho)\)-Poisson summation formula for split classical groups, using the standard representation of the complex dual group. The record omits the identity and its essential hypotheses, including the meaning of \(\sigma\), so this remains a statement-only draft requiring source clarification and human mathematical review. Its partially solved label is supplied dated curation.
\subsection{Sources}
\begin{itemize}
\item Supplied local UnsolvedMath registry record, permanent ID \(30004785\), title ``Direct Poisson Summation for Split Classical Groups,'' exact source problem code \texttt{catalogue entry}, in the Oberwolfach Reports Open Problems collection. The proposed year is recorded as 2021; no proposer is named.
\item Supplied source citation: \emph{Automorphic Forms, Geometry and Arithmetic (2022)}, Oberwolfach Reports, DOI 10.4171/owr/2021/39, pp.\ 2089--2156. \url{https://doi.org/10.4171/owr/2021/39}. The background describes Oberwolfach Reports as documenting mathematical workshops and their open-problem discussions. This citation is supplied provenance; its contents have not been independently consulted here.
\item Supplied literature source: ``Oberwolfach Report 39/2021.'' \url{https://doi.org/10.4171/owr/2021/39}.
\item Supplied literature source: ``A refined Poisson summation formula.'' \url{https://arxiv.org/abs/1707.06091}. The record supplies no authors or theorem locator for the claimed results.
\item Literature assessment checked at \texttt{2026-08-22}: supplied status \texttt{partially\_solved}. The supplied assessment says that the formula is proved when the relevant global functorial transfer to \(\operatorname{GL}_n\) is known and cuspidal, and that \(\operatorname{GL}_n\) has a direct theorem, while this conditional mechanism does not establish the conjecture for every pair \((G,\rho)\). This is supplied dated curation derived from the OWR material, not an independently verified literature assessment. The record expressly requests independent verification of time-sensitive status and cited claims.
\end{itemize}
\end{document}
